\section*{Chapitre \ref{ch:g8:balanced-equations} --- Conservation de la masse et équations ajustées}

\begin{solution}{exo:g8:balanced-equations:1}
Au cours d'une \omterm{def:g7:chemical-reactions:reaction}{réaction chimique}, la masse totale des produits formés
est égale à la masse totale des \omterm{def:g7:chemical-reactions:reactant}{réactifs} consommés.
\end{solution}

\begin{solution}{exo:g8:balanced-equations:2}
Le coefficient est 3. Il représente 3 \omterm{def:g7:atoms-and-molecules:atom}{atomes} de carbone et
$3 \times 2 = 6$ \omterm{def:g7:atoms-and-molecules:atom}{atomes} d'oxygène.
\end{solution}

\begin{solution}{exo:g8:balanced-equations:3}
Oui~: 1 C et 2 O de chaque côté.
\end{solution}

\begin{solution}{exo:g8:balanced-equations:4}
\ce{2Mg + O2 -> 2MgO}~: 2 Mg et 2 O de chaque côté.
\end{solution}

\begin{solution}{exo:g8:balanced-equations:5}
Azote~: 2 N à gauche, donc 2 \ce{NH3}~; il y a alors 6 H à droite, donc
3 \ce{H2}~: \ce{N2 + 3H2 -> 2NH3}.
\end{solution}

\begin{solution}{exo:g8:balanced-equations:6}
Par \omterm{def:g8:balanced-equations:conservation}{conservation de la masse}~: $44 - 12 = \qty{32}{g}$ de dioxygène.
\end{solution}

\begin{solution}{exo:g8:balanced-equations:7}
Le gaz formé s'échappe dans la pièce, et la balance ne le pèse plus.
Aucune masse n'est détruite~: la masse manquante est dans l'\omterm{def:g4:air-a-mixture-of-gases:air}{air} de la
pièce.
\end{solution}

\begin{solution}{exo:g8:balanced-equations:8}
Carbone~: 1 de chaque côté. Hydrogène~: 4 à gauche, donc 2 \ce{H2O}.
Oxygène~: $2 + 2 = 4$ à droite, donc 2 \ce{O2}~:
\ce{CH4 + 2O2 -> CO2 + 2H2O}.
\end{solution}

\begin{solution}{exo:g8:balanced-equations:9}
Changer \ce{H2O} en \ce{H2O2} change la substance~: \ce{H2O2} est le
peroxyde d'hydrogène, pas l'eau. Seuls les coefficients peuvent être
modifiés~: \ce{2H2 + O2 -> 2H2O}.
\end{solution}

\begin{solution}{exo:g8:balanced-equations:10}
$20 \times 5 = 100$ \omterm{def:g7:atoms-and-molecules:molecule}{molécules} de dioxygène~; $20 \times 4 = 80$
\omterm{def:g7:atoms-and-molecules:molecule}{molécules} d'eau.
\end{solution}

\begin{solution}{exo:g8:balanced-equations:11}
Avec 2 \ce{C2H6}~: 4 C, donc 4 \ce{CO2}~; 12 H, donc 6 \ce{H2O}~;
oxygène à droite $8 + 6 = 14$, donc 7 \ce{O2}~:
\ce{2C2H6 + 7O2 -> 4CO2 + 6H2O}.
\end{solution}

\begin{solution}{exo:g8:balanced-equations:12}
Toujours \qty{850.0}{g}. La boîte est scellée~: les \qty{0.4}{g} de
cire qui ont brûlé sont devenus du dioxyde de carbone et de la vapeur
d'eau, qui sont toujours dans la boîte avec le dioxygène qu'ils ont
consommé. Rien n'est entré ni sorti.
\end{solution}

\begin{solution}{pb:g8:balanced-equations:1}
\textbf{1.} Du dioxygène de l'\omterm{def:g4:air-a-mixture-of-gases:air}{air}, qui se combine au fer~: l'oxyde de
fer pèse le fer plus cet oxygène.

\textbf{2.} Tout reste dans le bocal~: le dioxygène consommé était déjà
dans le bocal, et l'oxyde formé y reste. La masse totale ne change pas.

\textbf{3.} La \omterm{def:g8:balanced-equations:conservation}{conservation de la masse}~: au cours d'une réaction, la
masse totale des produits est égale à la masse totale des \omterm{def:g7:chemical-reactions:reactant}{réactifs}.

\textbf{4.} \ce{4Fe + 3O2 -> 2Fe2O3}.

\textbf{5.} \ce{3Fe + 2O2 -> Fe3O4}.

\textbf{6.} À gauche~: 4 Fe, $3 \times 2 = 6$ O. À droite~:
$2 \times 2 = 4$ Fe, $2 \times 3 = 6$ O. L'équation est ajustée.

\textbf{7.} 4 \omterm{def:g7:atoms-and-molecules:atom}{atomes} de fer pour 3 \omterm{def:g7:atoms-and-molecules:molecule}{molécules} de dioxygène, donc $400 \times
\frac{3}{4} = 300$ \omterm{def:g7:atoms-and-molecules:molecule}{molécules} de dioxygène.

\textbf{8.} $0.28 \times \frac{3}{7} = \qty{0.12}{g}$ d'oxygène.

\textbf{9.} $0.28 + 0.12 = \qty{0.40}{g}$ d'oxyde de fer.

\textbf{10.} Oui~: il en faut \qty{0.12}{g} et le bocal en contenait
environ \qty{0.5}{g}.

\textbf{11.} Le dioxygène consommé a quitté l'\omterm{def:g4:air-a-mixture-of-gases:air}{air} du bocal (il est
maintenant dans l'oxyde solide)~: le bocal contient moins de gaz
qu'avant, et l'\omterm{def:g4:air-a-mixture-of-gases:air}{air} de la pièce s'y engouffre. L'indication de la balance
augmente alors, de la masse de l'\omterm{def:g4:air-a-mixture-of-gases:air}{air} qui entre.

\textbf{12.} \qty{0.12}{g} de dioxygène ont été pris à l'\omterm{def:g4:air-a-mixture-of-gases:air}{air} du bocal.
\end{solution}
